\structuralelement{Введение}

\subsection*{Актуальность темы исследования}
\addcontentsline{toc}{subsection}{Актуальность темы исследования}

В последние годы сервисные мобильные роботы быстро проникают в бытовую
сферу: от роботов-пылесосов и автоматических курьеров до персональных
ассистентов, сопровождающих пользователя в домашних условиях. Ключевой
сценарий применения таких систем~--- следование за человеком в
ограниченной среде, где одновременно должны обеспечиваться безопасность
движения, устойчивость поведения при временной потере наблюдаемости
цели и учёт динамически меняющейся обстановки (перемещаемые предметы
быта, домашние животные, другие люди в квартире).

Классические подходы к управлению мобильными роботами, основанные на
раздельных контурах глобального и локального планирования и простых
реактивных алгоритмах типа Dynamic Window Approach, плохо справляются
с одновременной формализацией нескольких разнородных требований~---
удержание дистанции и центрирование цели, избежание столкновений,
сохранение плавности управления. Современное развитие вычислительной
техники и стохастических методов оптимизации сделало практически
применимыми методы модельного прогнозирующего управления (MPC), в том
числе их стохастическую выборка-ориентированную реализацию MPPI,
поддерживающую нелинейные модели и негладкие составные функции потерь.

Применение MPPI в задаче сопровождения пользователя в квартире
непосредственно отражает современные тренды и является актуальной
инженерной и научно-прикладной задачей.

\subsection*{Исследованность темы}
\addcontentsline{toc}{subsection}{Исследованность темы}

Кинематические и динамические модели омни-роботов с колёсами Mecanum
детально исследованы в работах по мобильной робототехнике; в литературе
известны соотношения, связывающие скорости вращения колёс с
обобщёнными скоростями корпуса и учитывающие ограничения исполнительных
механизмов~\cite{MorenoCaireta2020MPC}. Применение MPC для
Mecanum-роботов в задачах обхода препятствий рассматривалось в ряде
работ, демонстрирующих преимущества оптимизационного подхода по сравнению
с классическими реактивными регуляторами~\cite{MorenoCaireta2020MPC,Alkhaddad2023MPC}.

Стохастическая реализация MPC --- MPPI --- получила практическое
применение в навигационном стеке ROS~2 (Nav2), где имеется полноценная
реализация соответствующего локального планировщика. Параллельно широко
развиваются алгоритмы SLAM (Gmapping, Hector~SLAM,
Cartographer, slam\_toolbox), позволяющие одновременно решать задачу
построения карты и локализации робота в помещении~\cite{Sreerambatla2021SLAM,Nguyen2022SLAM}.

Вместе с тем целостное решение, объединяющее в одной воспроизводимой
среде задачу сопровождения пользователя (трекинг цели по камере +
позиционирование в карте), MPPI-контроллер с критиками, настроенными
под прикладные требования (дистанция до цели, центрирование, обход
динамических препятствий, устойчивое поведение при потере цели), в
литературе практически не описано в открытом виде. Данная работа
отвечает на этот запрос.

\subsection*{Решаемая проблема}
\addcontentsline{toc}{subsection}{Решаемая проблема}

В работе решается проблема разработки программно-алгоритмического
комплекса автономного мобильного омни-робота, который в условиях
типовой квартирной среды с подвижными препятствиями способен
удерживать подвижного пользователя в рабочей зоне камеры и
поддерживать заданную дистанцию, избегая столкновений и сохраняя
устойчивое поведение при временной потере цели.

\subsection*{Цель работы}
\addcontentsline{toc}{subsection}{Цель работы}

Целью ВКР является разработка и экспериментальное исследование
программного обеспечения автономного мобильного омни-робота,
обеспечивающего сопровождение пользователя в условиях квартирной среды
с динамическими препятствиями на базе стохастической реализации
модельного прогнозирующего управления.

\subsection*{Задачи работы}
\addcontentsline{toc}{subsection}{Задачи работы}

Для достижения поставленной цели сформулированы следующие задачи:
\begin{enumerate}
    \item Проанализировать существующие подходы к локальному планированию
          траектории мобильных роботов в ограниченной среде, в том числе
          методы модельного прогнозирующего управления и их стохастические
          реализации.
    \item Сформировать математическую модель омни-платформы с колёсами
          Mecanum, включая кинематические соотношения и ограничения
          управления.
    \item Формализовать составную функцию потерь, отражающую прикладные
          требования.
    \item Разработать воспроизводимую программно-симуляционную среду на
          базе Docker, ROS~2 и Gazebo для отладки и валидации алгоритмов
          управления.
    \item Реализовать модель омни-робота с 2D-лидаром и фронтальной
          RGB-камерой, создать узлы детекции цели и оценки ее положения.
    \item Реализовать критик MPPI-контроллера,
          обеспечивающий удержание цели в центре кадра камеры, и
          интегрировать его в навигационный стек Nav2.
    \item Реализовать критик MPPI-контроллера, формирующий
          направленный (анизотропный) штраф на основе локального
          поля препятствий, для повышения безопасности голономной
          омни-платформы при сближении с препятствиями.
    \item Автоматизировать процесс сбора статистики и визуализации результатов.
\end{enumerate}

\subsection*{Практическая значимость работы}
\addcontentsline{toc}{subsection}{Практическая значимость работы}

Результаты работы могут быть использованы:
\begin{itemize}
    \item как основа прототипа персонального сервисного робота,
          выполняющего автономное сопровождение пользователя в бытовой
          среде;
    \item как воспроизводимый программно-симуляционный стенд для
          исследования и настройки алгоритмов локального планирования и
          трекинга в условиях квартиры;
    \item для сравнения различных конфигураций критиков MPPI и
          подтверждения того, что декомпозиция составной функции потерь
          на независимые критики является практичным способом формализации
          разнородных требований к поведению робота.
\end{itemize}